\section{Threshold exteriors and geometric barriers}\label{four:sec:geometry}

We work with the prepared one-ended data furnished by
Lemma~\ref{four:lem:strictification}.  Thus the original boundary has strictly
negative future expansion, $K$ has compact support, and, for a smooth positive
extension $r$ of the end radius,
\begin{equation}\label{four:geo:prepared-gap}
 \mu-|J|_g\ge c_0r^{-4-\delta_1},\qquad c_0>0,\qquad 0<\delta_1<1.
\end{equation}
The enclosing infimum for this prepared metric is denoted by $A_*$, with the
intrinsic-boundary convention of Definition~\ref{four:def:cuts}.  This section
constructs an exterior on which both signs of a small prescribed trace can be
controlled.  The original trapping convention is always the future convention
$H+\tr_S K<0$.  The tensor $-K$ below is used only to construct additional
interior barriers.

\subsection{Total regions inside a compact barrier}

For a smooth symmetric tensor $Q$ and an oriented hypersurface $\Sigma$, write
\[
 \Theta_Q(\Sigma)=H_\Sigma+\tr_\Sigma Q.
\]
When $\Sigma$ bounds a compact region, its normal points out of that region.
A smooth compact region may have several components.  It includes its entire
manifold boundary and is the closure of its interior.  For a compact manifold
$X$ with boundary and a number $c$, define
\begin{equation}\label{four:geo:total-definition}
 T_c(X,Q)=\overline{\bigcup\left\{E:
 E\Subset\operatorname{int}X\text{ is a smooth compact region},\quad
 \Theta_Q(\partial E)\le c\right\}}.
\end{equation}
The empty union has empty closure.  All closures in this definition are taken
in $X$.

\begin{lemma}[Compact total region]\label{four:geo:compact-total}
Let $(X^4,g)$ be smooth, compact, connected, and orientable, with nonempty
smooth boundary, and let $Q$ be smooth.  Suppose that
$\Theta_Q(\partial X)>c$, with the normal pointing out of $X$ on every
boundary component.  Then $T_c(X,Q)$ is either empty or a smooth compact
region contained in $\operatorname{int}X$.  In the latter case its frontier
has expansion $c$ and is stable for outward variations.  It contains every
region in \eqref{four:geo:total-definition} and is itself one of those regions.
In particular, no component of its complement whose closure is contained in
$\operatorname{int}X$ can be a bounded hole.
\end{lemma}

\begin{proof}
Replace $Q$ by $Q_c=Q-(c/3)g$.  Since the hypersurfaces have dimension three,
\begin{equation}\label{four:geo:threshold-shift}
 \Theta_{Q_c}(\Sigma)=\Theta_Q(\Sigma)-c
\end{equation}
for every oriented hypersurface.  It suffices to prove the result for
$c=0$, and we use $Q_c$ throughout this proof.

The global input is the smooth-frontier part of
\cite[Theorem~4.6]{AnderssonEichmairMetzger2011}: for smooth complete asymptotically flat data in
spatial dimensions $2\le n\le7$, a nonempty total trapped region has smooth
embedded outermost stable MOTS frontier.  No dominant energy condition is
required for this conclusion.  We explain how to meet its boundaryless
asymptotically flat hypotheses while preserving exactly the compact problem.

First attach an outward cobordism from $\partial X$ to one three-sphere.
Each component of $\partial X$ is a closed orientable three-manifold.  Relative
one-handles join these components, and the smooth surgery theorem for
orientable three-manifolds gives a framed link surgery from the resulting
connected three-manifold to $S^3$ \cite[Theorem~6, p.~522]{Wallace1960}.  Its trace consists of
relative two-handles: reversing a four-dimensional two-handle trace again
uses two-handles.  Thus the cobordism has a handle decomposition relative to
$\partial X$ with indices at most two.  Attach $S^3\times[0,\infty)$ at its
other end, obtaining a boundaryless manifold $\widetilde X$.  The handle
Morse function provides a proper smooth height $q$
on the extension, with $q=0$ on $\partial X$, increasing in an outward
collar, with only critical points of index at most two, and equal to a radial
coordinate sufficiently far out.  Extend $q$ a short distance into $X$ as a
collar coordinate with negative values.

Extend $g$ smoothly across $\partial X$.  We may choose its values near every
critical point of $q$ so that
\begin{equation}\label{four:geo:morse-trace}
 \tr_V\Hess_g q>0
 \quad\text{for every three-dimensional subspace }V\subset T_x\widetilde X
\end{equation}
at that critical point, and hence throughout a smaller neighborhood.  Indeed,
at an index-two point the eigenvalues of the Hessian can be arranged as
$-a,-b,A,B$, with $a,b>0$ and $\min(A,B)>a+b$; the sum of the three smallest
is positive.  At indices zero or one the same choice is easier.  This is
achieved by rescaling the positive and negative coordinate directions in
the metric in a Morse chart.  The neighborhoods are disjoint and away from
the prescribed original data, so these local metrics extend to a smooth
positive metric on the whole extension.  Choose that metric Euclidean on
the remote end.

At a regular point of a level of $q$, its mean curvature toward increasing
$q$ is
\[
 H_{\{q=\text{constant}\}}
 =\frac{\tr_{(\grad q)^\perp}\Hess q}{|\dd q|}.
\]
It is positive near the critical points by \eqref{four:geo:morse-trace}.
Choose a smooth preliminary tensor extension agreeing with $Q_c$ in the
prescribed boundary collar. Near each critical point there is $a>0$ such
that $\tr_V\Hess q\ge a$ on every unit three-plane. After shrinking its
neighborhood, $H\ge a/|\dd q|$ dominates the bounded tangential trace of
the preliminary tensor. On the remaining compact regular part the level
mean curvature is bounded below. Retain the tensor near the original
boundary, and add a sufficiently large nonnegative multiple of the metric
on the remaining compact part of the extension.  The tangential trace of this addition is three times its
coefficient.  We can therefore ensure strictly positive expansion on every
regular level of $q$.  The given strict boundary inequality permits the
matching in a collar of $\partial X$, including a collar on the original
side.  On the Euclidean remote end take the tensor to be zero; its round
levels already have positive mean curvature.  The transition to zero can
be made inside that Euclidean region with a nonnegative metric multiple.
We have obtained smooth complete one-ended data without boundary, exactly
Euclidean with zero tensor sufficiently far out.

Every compact smooth weakly trapped region in this extension stays a
positive distance inside $X$.  To see this, extend $q$ to a smooth function
on the interior core with values below those of a fixed smaller collar.
If a region enters that collar or the attached extension, the maximum of
$q$ on it occurs at a point in that part.  A regular maximum cannot lie in
the interior of the region.  At a regular boundary maximum the region is
on the inner side of a positively expanding level, with the same outward
normal at contact.  Tangential Hessian comparison gives
$\Theta_{Q_c}(\partial E)\ge\Theta_{Q_c}(\{q=q_{\max}\})>0$, a
contradiction.  At a critical maximum in the interior, a direction of
positive Hessian contradicts maximality.  At a critical boundary maximum,
the inward open half-space of a smooth domain contains a vector $v$ with
$\Hess q(v,v)>0$: choose the sign of a positive Hessian direction to point
inward, and perturb it inward if it is tangent.  A curve entering the region
with this initial velocity has
$q(\gamma(t))=q(\gamma(0))+t^2\Hess q(v,v)/2+o(t^2)$, again a
contradiction.  Thus one fixed original-side collar is avoided by every
such region.

The global total region consequently coincides with
\eqref{four:geo:total-definition}; all its candidate regions lie in the same
compact subset of $\operatorname{int}X$.  The cited theorem gives its smooth
stable frontier.  With the closed representative used here it is a smooth
compact region, its frontier has $\Theta_{Q_c}=0$, and it contains every
candidate.  This proves the assertion at threshold $c$ by
\eqref{four:geo:threshold-shift}.  If a complementary component has closure in
$\operatorname{int}X$, filling it deletes whole smooth boundary components
and changes no remaining outward expansion.  The resulting larger region
would still be a candidate in \eqref{four:geo:total-definition}, contradicting
maximality.
\end{proof}

The extension in this proof is auxiliary.  Neither a constraint inequality
nor a topological restriction is imposed on it.  We will also apply the
lemma after compactly supported changes of $Q$ that leave the strict outer
barriers unchanged.

\subsection{Threshold continuity and collars}

\begin{lemma}[Continuity thresholds]\label{four:geo:threshold-continuity}
Suppose that the strict barrier condition in Lemma~\ref{four:geo:compact-total}
holds for every $c$ in an open interval $I$.  The closed regions $T_c(X,Q)$
increase with $c$.  Except at a countable set of thresholds, they are
right-continuous in Hausdorff distance.  At an empty right-continuity
threshold the regions are empty for all sufficiently close larger
thresholds.
\end{lemma}

\begin{proof}
Inclusion of the defining classes proves monotonicity.  Let $L$ exceed the
diameter of $X$, and associate to each threshold the function
\[
 d_c(x)=\begin{cases}
 \operatorname{dist}(x,T_c(X,Q)),&T_c(X,Q)\ne\varnothing,\\
 L,&T_c(X,Q)=\varnothing.
 \end{cases}
\]
These functions are uniformly bounded, are $1$-Lipschitz in $x$, and are
nonincreasing in $c$.  Fix a countable dense set $\{x_j\}$ in $X$.  For each
$j$ the monotone real function $c\mapsto d_c(x_j)$ has at most countably
many discontinuities.  Away from the union of these exceptional sets,
$d_{c_i}(x_j)\to d_c(x_j)$ whenever $c_i\downarrow c$.  A finite net in $X$
and the common Lipschitz bound promote convergence on the dense set to
uniform convergence on $X$.  For two nonempty compact sets the uniform
distance between their distance functions equals their Hausdorff distance.
If $T_c$ is empty, uniform convergence to the constant $L>\operatorname{diam}X$
is impossible along nonempty $T_{c_i}$, whose distance functions vanish
somewhere.  Thus eventual emptiness also follows.  In particular every
nonempty interval contains an admissible right-continuity threshold.
\end{proof}

\begin{lemma}[Outward leaf collars]\label{four:geo:stable-collar}
Let $T_c(X,Q)$ be nonempty, and let $\Sigma$ be one connected component of
its frontier.  On the side exterior to this region, $\Sigma$ has a smooth
foliated collar, parametrized by $s\in[0,s_0]$, with $|\dd s|>0$.  With
normal $N=\grad s/|\dd s|$, every leaf satisfies
\begin{equation}\label{four:geo:collar-expansion}
 \Theta_Q(\Sigma_s)\ge c.
\end{equation}
The collars of distinct components can be chosen disjoint.
\end{lemma}

\begin{proof}
Write $\Theta(v)$ for the expansion of the outward normal graph of a small
function $v$ over $\Sigma$, pulled back to $\Sigma$.  If $\nu$ is the
outward normal and $\mathrm{II}$ its second fundamental form, the normal
variation formulas give its linearization
\[
 Lv=-\Delta_\Sigma v+2Q(\nu,\grad_\Sigma v)
   +\bigl(\tr_\Sigma\nabla_\nu Q-|\mathrm{II}|^2
                          -\Ric(\nu,\nu)\bigr)v.
\]
Thus $L$ is a smooth scalar elliptic operator with principal part
$-\Delta_\Sigma$.  Stability gives a nonnegative principal eigenvalue
$\lambda_1$, with a positive eigenfunction $\phi$.  The principal
eigenvalue of the adjoint is the same and has a positive eigenfunction
$\phi^*$; these spectral facts apply on a compact connected manifold of
any dimension \cite[Lemma~4.1]{AnderssonMarsSimon2008}.  The shift in
\eqref{four:geo:threshold-shift} has zero derivative under surface variations,
so this is also the stability operator for expansion $\Theta_Q=c$.

If $\lambda_1>0$, the graphs $v=s\phi$ have
$\Theta(v)=c+s\lambda_1\phi+O(s^2)>c$ for all sufficiently small $s>0$.
Their positive normal speed makes them a collar.  It remains to treat
$\lambda_1=0$.

Fix $0<\alpha<1$.  The bounded linear map
\begin{equation}\label{four:geo:augmented-stability}
 C^{2,\alpha}(\Sigma)\times\R\longrightarrow C^\alpha(\Sigma)\times\R,
 \qquad (v,a)\longmapsto
 \left(Lv-a,\int_\Sigma v\,\dd A\right)
\end{equation}
is an isomorphism.  Indeed, for prescribed $(f,b)$, pairing $Lv-a=f$
with $\phi^*$ determines
\[
 a=-\frac{\int_\Sigma\phi^*f\,\dd A}
          {\int_\Sigma\phi^*\,\dd A}.
\]
The resulting right side $f+a$ is orthogonal to the adjoint kernel, so the
Fredholm alternative solves $Lv=f+a$.  The kernel is the span of $\phi$;
adding its unique multiple that gives integral $b$ proves existence and
uniqueness in \eqref{four:geo:augmented-stability}.  Elliptic estimates give a
bounded inverse.

The implicit-function theorem applied to
\[
 (v,a)\longmapsto
 \left(\Theta(v)-c-a,\int_\Sigma v\,\dd A\right)
\]
therefore produces smooth graphs $v_s$ of constant expansion $c+a_s$ with
$v_0=0$, $a_0=0$, and $\int_\Sigma v_s\,\dd A=s$.  Differentiation at zero
and pairing with $\phi^*$ give
\[
 a'_0=0,\qquad v'_0=\frac{\phi}{\int_\Sigma\phi\,\dd A}>0.
\]
After shortening the parameter interval, the graphs have positive normal
speed and lie strictly outside $T_c$ for $s>0$.  If $a_s\le0$ for any such
$s$, replace $\Sigma$ by this outward graph and leave the other boundary
components unchanged.  The resulting smooth compact region would have
expansion at most $c$ and would strictly enlarge $T_c$.  This contradicts
Lemma~\ref{four:geo:compact-total}.  Thus $a_s>0$ for every sufficiently small
$s>0$, proving \eqref{four:geo:collar-expansion} also in the degenerate case.
Smooth elliptic regularity gives smooth leaves and parameter dependence.
Compactness of the frontier gives finitely many components; sufficiently
short collars are mutually disjoint.
\end{proof}

\subsection{Exterior supports and smooth enclosing regions}

A \emph{smooth exterior support} to a closed set $D$ at
$x\in\partial D$ is a smooth regular level $\{\psi=0\}$ through $x$,
defined in a neighborhood of $x$, such that $D$ is locally contained in
$\{\psi\le0\}$.  Its outward normal is $\grad\psi/|\dd\psi|$.  Saying
that $D$ has expansion at most $k$ in exterior supports means that every
such support has $\Theta_Q\le k$ at contact.  This definition makes no
regularity assumption on $D$.

The next argument is useful because the curvatures of these supports need
not be bounded.  Ordinary approximate transport of tangent metrics would
produce an uncontrolled error multiplying their second fundamental forms.
We use a map whose differential is an exact isometry at the contact point.

\begin{lemma}[Exact support transport]\label{four:geo:support-transport}
Let $g,Q$ be smooth on a neighborhood of a compact set, and let $D$ be a
compact closed subset of its interior.  Suppose that $D$ has expansion at
most $k$ in exterior supports.  There are $z_0>0$ and $C<\infty$ such that,
for $0<z<z_0$, every exterior support to
\[
 D_z=\{y:\operatorname{dist}_g(y,D)\le z\}
\]
has expansion at most $k+Cz$.  The constant $C$ depends only on bounds for
the smooth background geometry and $Q$ in the fixed compact neighborhood,
and not on the support or its curvature.  The allowable $z_0$ is also
smaller than the distance required to keep these neighborhoods in the
interior.
\end{lemma}

\begin{proof}
The assertion is immediate if $D$ is empty.  Let a support $\{\psi=0\}$
touch $D_z$ at $x'$, and choose $x\in D$ with
$\operatorname{dist}(x,x')=z$.  This nearest point lies in $\partial D$,
since an interior point could be moved a short distance toward $x'$ to
decrease the distance.  Choose $z_0$ below a uniform injectivity
radius.  Let $\gamma:[0,z]\to X$ be the unit-speed minimizing segment from
$x$ to $x'$, put $u_0=\dot\gamma(0)$, and denote parallel transport along
$\gamma$ by $P_t$.  The ball of radius $z$ centered at $x$ lies in $D_z$.
Consequently the support normal at $x'$ is $P_z u_0$.

For each $X\in T_xX$, solve the Jacobi boundary problem
\begin{equation}\label{four:geo:jacobi-transport}
 D_t^2J_X+R(J_X,\dot\gamma)\dot\gamma=0,
 \qquad J_X(0)=X,\qquad J_X(z)=P_zX.
\end{equation}
Here $R$ denotes the Riemann curvature tensor.  This problem is uniquely
solvable for uniformly small $z$.  To see the estimates explicitly, write
$J_X(t)=P_t(X+W_X(t))$.  The equation for $W_X$ has zero endpoint values
and is
$W_X''=-\mathcal R_t(X+W_X)$, with bounded curvature endomorphism
$\mathcal R_t$.  The Dirichlet inverse of the second derivative on
$[0,z]$ has supremum-norm bound $Cz^2$.  Absorption for small $z$, followed
by the differentiated one-dimensional Green formula, gives
\begin{equation}\label{four:geo:jacobi-jet-bounds}
 \sup_{[0,z]}|W_X|\le Cz^2|X|,
 \qquad |D_tJ_X(0)|\le Cz|X|.
\end{equation}
The longitudinal component of $J_X$ is affine with equal endpoint values.
Hence the linear map $A_zX=D_tJ_X(0)$ satisfies
\[
 \inner{A_zX}{u_0}=0,\qquad |A_z|\le Cz.
\]
These are precisely the first-jet compatibility conditions for a unit
vector field $U$ with $U(x)=u_0$ and $\nabla U(x)=A_z$.

For completeness, use normal coordinates $y^i$ at $x$ and extend vectors
by parallel transport along the radial geodesics of those coordinates.
Writing $E_v(y)$ for the resulting extension of $v\in T_xX$, set
\[
 V(y)=E_{u_0}(y)+\sum_i y^i E_{A_ze_i}(y),\qquad
 U(y)=\frac{V(y)}{|V(y)|}.
\]
On a uniformly small neighborhood $V$ stays away from zero.  At $x$ this
field has the required value and first derivative, since $A_zX\perp u_0$.
It has uniformly bounded first and second derivatives by smoothness and
compactness of the background.  No vanishing second covariant derivative
is required.

Define the local map
\[
 F_z(y)=\exp_y(zU(y)).
\]
The geodesic variation determining $\dd F_z(x)X$ has initial Jacobi data
$X,A_zX$, so \eqref{four:geo:jacobi-transport} gives the exact identities
\begin{equation}\label{four:geo:exact-isometry}
 F_z(x)=x',\qquad \dd F_z(x)=P_z,
 \qquad |\nabla\dd F_z(x)|\le Cz.
\end{equation}
To verify the last estimate, the smooth map $E(y,v)=\exp_y(v)$ satisfies
$E(y,0)=y$.  Differentiating $E(y,zU(y))$ twice in any fixed coordinate
chart, with the uniform $C^2$ bound on $U$, gives
$D F_z=I+O(z)$ and $D^2F_z=O(z)$.  In the covariant second derivative the
domain and target connection terms cancel at $z=0$; their remaining
difference is $O(z)$.  This proves the estimate uniformly in $x,x'$ and
$u_0$.

For $y\in D$ sufficiently close to $x$, the curve defining $F_z(y)$ has
length $z$, and therefore $F_z(y)\in D_z$.  Thus $\psi\circ F_z$ is an
exterior support to $D$ at $x$.  Its gradient is nonzero by the exact
isometry in \eqref{four:geo:exact-isometry}, and its normal is $u_0$.  For an
orthonormal basis $e_1,e_2,e_3$ of $u_0^\perp$, the Hessian chain rule is
\[
 \Hess(\psi\circ F_z)(e_a,e_a)
 =\Hess\psi(P_ze_a,P_ze_a)
   +\dd\psi\bigl((\nabla\dd F_z)(e_a,e_a)\bigr).
\]
Since $|\dd(\psi\circ F_z)|_x=|\dd\psi|_{x'}$, taking the tangential
trace and dividing by this common norm yields
\begin{equation}\label{four:geo:mean-curvature-transport}
 H_{\psi\circ F_z}(x)-H_\psi(x')
 =\sum_{a=1}^3
   \inner{P_zu_0}{(\nabla\dd F_z)(e_a,e_a)}=O(z).
\end{equation}
There is no error involving $\Hess\psi$ in this formula.  Smoothness of
$Q$ similarly gives
\[
 \left|\tr_{u_0^\perp}Q(x)
       -\tr_{P_z(u_0^\perp)}Q(x')\right|\le Cz.
\]
Apply the assumed support inequality at $x$, then combine these two bounds.
The expansion of the original support at $x'$ is at most $k+Cz$, as claimed.
\end{proof}

\begin{lemma}[Smooth support enclosure]\label{four:lem:support-enclosure}
Assume the hypotheses of Lemma~\ref{four:geo:compact-total} at threshold $c'$.
Let $D\Subset\operatorname{int}X$ be a compact closed set whose every
smooth exterior support has $\Theta_Q\le k$, where $k<c'$.  Then
\[
 D\subset T_{c'}(X,Q).
\]
In particular, the right side is nonempty whenever $D$ is nonempty.
\end{lemma}

\begin{proof}
Assume $D\ne\varnothing$.  Choose $\beta>0$ small enough that $D_\beta$
lies in $\operatorname{int}X$, Lemma~\ref{four:geo:support-transport} applies
for $0<z\le\beta$, and
\begin{equation}\label{four:geo:strict-support-gap}
 k+C\beta<c'.
\end{equation}
Smooth approximation of the distance function, followed by a regular-value
choice, gives a smooth compact region $E$ with
\[
 D_{\beta/4}\subset\operatorname{int}E
 \subset E\subset\operatorname{int}D_{3\beta/4}.
\]
For example, approximate the distance uniformly to accuracy $\beta/16$
and choose a regular level between $7\beta/16$ and $9\beta/16$.

Choose a smooth function $a\ge0$, supported where
$\operatorname{dist}(\cdot,D)<\beta$, that is a sufficiently large
constant near $\partial E$.  For $Q'=Q-ag$ we then have
$\Theta_{Q'}(\partial E)=\Theta_Q(\partial E)-3a<c'$.
The boundary barriers of $X$ are unchanged.  Lemma~\ref{four:geo:compact-total}
applied to $Q'$ produces a nonempty smooth total region $T'$ containing
$E$.  Its frontier is nonempty because $T'$ is compactly contained in the
interior of the connected manifold $X$ with nonempty boundary.  That
outward frontier satisfies
\begin{equation}\label{four:geo:modified-frontier}
 \Theta_Q(\partial T')=c'+3a\ge c'.
\end{equation}
Since $D$ has an open neighborhood in $E\subset T'$, the number
$z=\operatorname{dist}(D,\partial T')$ is positive and attained.

Suppose $z<\beta$.  Then $D_z\subset T'$.  Indeed, a minimizing segment
of length at most $z$ from $D$ to a point outside $T'$ would first cross
$\partial T'$ at distance strictly less than $z$.  Such segments stay in
$X$ by our choice of $\beta$.  At a pair realizing the minimum distance,
$\partial T'$ is therefore an exterior support to $D_z$.  Its expansion
is at most $k+Cz<c'$ by Lemma~\ref{four:geo:support-transport} and
\eqref{four:geo:strict-support-gap}, contradicting
\eqref{four:geo:modified-frontier}.  Consequently $z\ge\beta$.

The frontier of $T'$ is outside the support of $a$, so its expansion for
the original tensor is exactly $c'$.  Thus $T'$ itself is an admissible
region in \eqref{four:geo:total-definition} for $(X,Q,c')$.  Since it contains
$D$, the defining maximality gives $D\subset T'\subset T_{c'}(X,Q)$.
\end{proof}

\subsection{The two colors and the trace allowance}

\begin{proposition}[Geometric preparation]\label{four:prop:geometric-preparation}
For the prepared data satisfying \eqref{four:geo:prepared-gap}, there are a
closed connected exterior $\Omega\subset M$ with one end and nonempty
compact smooth boundary $B$, two negative numbers $c_b,c_w$, and a
decomposition of $B$ into black and white components with the following
properties.  With the normal $\nu$ into $\Omega$ and $P_B=\tr_B K$,
\begin{equation}\label{four:geo:two-boundary-expansions}
 H+P_B=c_b\quad\text{on black components},\qquad
 H-P_B=c_w\quad\text{on white components}.
\end{equation}
The components come from closed total-region families $\cD_c$ and $\cW_c$
at right-continuity thresholds $c_b$ and $c_w$, respectively.  The white
family is formed with the black region fixed, and is allowed to be empty.
Every face has a disjoint smooth collar in $\Omega$, with parameter $s\ge0$,
$|\dd s|>0$, and the corresponding expansion toward increasing $s$ is
at least $c_b$ or $c_w$.

There are $a_0>0$, a compactly supported smooth function $C$, and a positive
smooth weight $\rho$, equal to a positive constant times
$r^{-4-\delta_1}$ on the end, such that $|C|\le a_0$ and, on smaller
collars,
\begin{equation}\label{four:geo:linear-offsets}
 C=a_0-\kappa_Bs\quad\text{on black collars},\qquad
 C=-a_0+\kappa_Bs\quad\text{on white collars},\qquad \kappa_B>0.
\end{equation}
Set
\begin{equation}\label{four:geo:height-endpoints}
 b_-=c_b-a_0<0< -c_w+a_0=b_+.
\end{equation}
Every real-valued function $h$ taking values in $[b_-,b_+]$ on
$\supp K\cap\Omega$ satisfies the following inequality on $\Omega$:
\begin{equation}\label{four:eq:trace-slack}
 |(h+C)\tau|+|\dd C|_g+\rho\le\mu-|J|_g,
 \qquad \tau=\tr_g K.
\end{equation}
There is a smooth compactly supported extension $b$ on $\Omega$ of the
boundary values $b_-$ on black faces and $b_+$ on white faces, constant on
smaller collars.  Finally, every smooth enclosing cut of $B$ in $\Omega$
has $g$-volume at least $A_*$.
\end{proposition}

\begin{proof}
Choose a distant sphere $S_{R_0}$ beyond $\supp K$, sufficiently large that
it and all farther coordinate spheres have positive mean curvature.
Their expansions for both $K$ and $-K$ are positive.  Let $X_0$ be the
original compact core between $S$ and $S_{R_0}$.  It is connected: paths
using the remote end can be replaced at its connected spherical cross
section.  Attach a smooth compact filling $F$ behind the entire original
boundary.  Such a filling can be constructed by taking an oppositely
oriented copy of $X_0$ and capping its spherical boundary by a four-ball.
Its remaining boundary identifies with $S$.  Extend $g$ and $K$ smoothly
across $S$ into this filling, keeping the metric positive.  Extension of
the collar jets followed by a partition of unity does this; no constraint
condition is required in $F$.  The resulting compact manifold $X$ has
only the connected outer boundary $S_{R_0}$.

For later parameter choices, fix now
\[
 \sigma_* =\min_{X_0}(\mu-|J|_g)>0,
 \qquad T_* =\max_{X_0}|\tau|.
\]
Choose $\gamma>0$ so small that
\begin{equation}\label{four:geo:threshold-budget}
 \gamma T_*\le\sigma_*/8
\end{equation}
and the original boundary has expansion less than $-2\gamma$.
By strictness and continuity, $F$ together with a fixed short exterior
collar of $S$ is a smooth compact region whose outward boundary has
expansion less than $-\gamma$.  For every $c\in(-\gamma,0)$ form
\[
 \cD_c=T_c(X,K).
\]
Here $K$ denotes its smooth extension into $F$.  Lemma~\ref{four:geo:compact-total}
applies and $\cD_c$ contains that entire filled neighborhood.  Its smooth
frontier therefore lies strictly in the original data.  Its complement
in $X$ is connected.  Indeed the collar of the connected outer sphere
belongs to a unique complementary component; every other component would
be an interior hole and could be filled, contrary to maximality.

Choose $c_b\in(-\gamma,0)$ at a right-continuity threshold by
Lemma~\ref{four:geo:threshold-continuity}.  Write $Y$ for the closed complement
of $\operatorname{int}\cD_{c_b}$ in $X$.  This is a smooth compact connected
manifold with the outer sphere and the black frontier as its boundary.
For the white tensor $-K$, all these boundary components have strictly
positive outward expansion.  On a black component the outward normal of
$Y$ is the reverse of the outward normal of $\cD_{c_b}$, so the expansion
is
\begin{equation}\label{four:geo:reversed-black-barrier}
 -H-\tr_{\partial\cD_{c_b}}K=-c_b>0.
\end{equation}
This changes both the normal and the tensor in an auxiliary problem;
the original boundary still has its prescribed future sign.

For $c\in(-\gamma,0)$ define the white family
\[
 \cW_c=T_c(Y,-K).
\]
It may be empty.  Lemmas~\ref{four:geo:compact-total} and
\ref{four:geo:threshold-continuity} give a right-continuity threshold
$c_w\in(-\gamma,0)$.  Every nonempty white region lies strictly inside
$Y$, so $\cD_{c_b}$ and $\cW_{c_w}$ are disjoint closed sets, a positive
distance apart.  Notice that the black threshold and manifold $Y$ are
fixed before making this second threshold choice.

In the filled one-ended manifold take the component of the open
complement of $\cD_{c_b}\cup\cW_{c_w}$ that reaches the end, and let
$\Omega$ be its closure.  Its boundary is the union of those entire
frontier components adjoining this exterior.  To justify the word
``entire,'' a smooth connected frontier component has a connected thin
collar on each side; the component of its exterior side is therefore
constant along the frontier.  Both frontiers are compact and smooth with
finitely many components, and they are disjoint.  Thus $\Omega$ is smooth,
closed, and connected, and its boundary $B$ is a union of complete smooth
components.  It has one end, the original end, and lies in the original
manifold because the black region contains the filling and a collar of
$S$.  Its boundary is nonempty: the component reaching infinity is
separated from that nonempty filled region.  The inherited metric is
complete up to $B$; boundary charts give local completeness at finite
limit points, and completeness of the original end prevents escape in
finite distance.

Call its retained $\cD_{c_b}$ faces black and its retained $\cW_{c_w}$
faces white.  The normals into $\Omega$ are the outward normals of these
regions.  Their threshold equations give
\eqref{four:geo:two-boundary-expansions}.  Figure~\ref{four:fig:two-frontiers}
illustrates these orientations; $\Omega_R$ there denotes the part of
$\Omega$ inside a distant coordinate sphere $S_R$.
\begin{figure}[tbp]
\centering
\begin{tikzpicture}[>=Latex,scale=.95,every node/.style={font=\small}]
  \draw[densely dashed,gray] (0,0) ellipse (4.4 and 2.05);
  \node[fill=white,inner sep=2pt] at (3.7,1.2) {$S_R$};
  \filldraw[fill=black!13,draw=black,thick] (-1.75,0) ellipse (1.12 and 1.08);
  \draw[densely dashed] (-1.92,-.02) ellipse (.48 and .42);
  \node at (-1.92,-.02) {$S$};
  \node at (-1.75,.72) {$\cD_{c_b}$};
  \filldraw[fill=white,draw=black,thick] (1.5,.65) ellipse (.78 and .53);
  \node at (1.5,.65) {$\cW_{c_w}$};
  \draw[->,thick] (-.63,0) -- (.18,0) node[midway,above] {$\nu_b$};
  \draw[->,thick] (1.5,.12) -- (1.5,-.65)
        node[midway,right] {$\nu_w$};
  \node at (.15,-1.05) {$\Omega_R$};
  \node at (-1.75,-1.52) {$H+P_B=c_b$};
  \node at (1.5,1.55) {$H-P_B=c_w$};
  \draw[->] (4.4,0) -- (5.2,0);
  \node[right] at (5.2,0) {end \(e\)};
\end{tikzpicture}
\caption{A schematic of the auxiliary exterior; one frontier component of
each kind is shown. The black region encloses the original boundary.
Both normals point into the retained exterior, but its black and white
faces have prescribed expansion for \(K\) and \(-K\), respectively.
The drawing represents incidence and orientations, not the topology
or dimension of the hypersurfaces.}
\label{four:fig:two-frontiers}
\end{figure}


Apply
Lemma~\ref{four:geo:stable-collar} to their respective total regions and shorten
the collars so they lie on the $\Omega$ side and are mutually disjoint.
This gives all the asserted leaf inequalities.  The support enclosure
Lemma~\ref{four:lem:support-enclosure} is available for the full families
$\cD_c$ and $\cW_c$ on $X$ and $Y$, respectively; in the latter case the
reversed black barriers remain fixed.

All these collars lie in the fixed original core $X_0$, so the lower
bound $\sigma_*$ was fixed before either threshold or any collar width.
For each collar choose a smooth profile $q_B(s)\in[0,1]$, supported in
that collar, with $q_B(s)=1-\alpha_Bs$ near its face and
$\alpha_B>0$.  Such a profile is obtained by multiplying this linear
function by a nonnegative cutoff supported before the linear function
can become negative.  With disjoint supports, define
\[
 C=a_0\sum_B\varepsilon_Bq_B(s),\qquad
 \varepsilon_B=\begin{cases}1,&B\text{ black},\\-1,&B\text{ white}.
 \end{cases}
\]
There is a finite constant $M_*$, now depending on the chosen collars and
profiles, such that $|C|\le a_0$ and $|\dd C|\le a_0M_*$.  Formula
\eqref{four:geo:linear-offsets} holds with $\kappa_B=a_0\alpha_B>0$.
Only at this point choose $a_0>0$ small enough that
\begin{equation}\label{four:geo:amplitude-budget}
 a_0(2T_*+M_*)\le\sigma_*/8.
\end{equation}
For $b_-\le h\le b_+$ we have $|h+C|\le\gamma+2a_0$, and hence on
$X_0\cap\Omega$
\begin{equation}\label{four:geo:compact-slack}
 |(h+C)\tau|+|\dd C|
 \le\gamma T_*+a_0(2T_*+M_*)\le\sigma_*/4.
\end{equation}
At points outside $\supp K$ the trace term vanishes, so this estimate
requires no restriction there on the value of $h$.

Finally choose a constant $\rho_*>0$ so small that
\[
 \rho_*\le c_0/4,\qquad
 \rho_*\max_{X_0}r^{-4-\delta_1}\le\sigma_*/4,
 \qquad \rho=\rho_* r^{-4-\delta_1}.
\]
On $X_0\cap\Omega$, \eqref{four:geo:compact-slack} plus this choice bounds the
left side of \eqref{four:eq:trace-slack} by $\sigma_*/2$.  Outside $X_0$,
$K=C=0$, and \eqref{four:geo:prepared-gap} gives the required inequality
directly.  This proves \eqref{four:eq:trace-slack}.  This order of choices uses
no lower bound on collar widths as $c_b,c_w$ approach zero: the thresholds
are chosen first, the collars next, and their amplitude last.  Ordinary
collar cutoffs also give the stated compactly supported smooth extension
of the constants $b_-$ and $b_+$.

If $D_e$ is an admissible closed connected outer domain for a cut in
$\Omega$, it is also closed in $M$, because $\Omega$ is closed in $M$.
It is a smooth codimension-zero submanifold of $M$, its interior lies in
$\operatorname{int}M$, and it contains the distant original end.  Its
entire intrinsic boundary, including every part coincident with $B$, is
therefore an admissible original enclosing cut.  Its area is at least
$A_*$.  This proves the final assertion without assuming the existence of
a minimizing cut.
\end{proof}
